% ============================================================
%  Chapter 1 — Abbreviations and Notations
% ============================================================
\section{Some Abbreviations and Notations}

These notes follow a set of standard shorthand abbreviations and mathematical
notations that are used consistently throughout.  Familiarity with them from the
outset will make the remaining chapters much easier to read.

% ------------------------------------------------------------
\begin{defbox}[Definition --- Abbreviations Used in These Notes]
The following abbreviations appear throughout:
\begin{center}
\renewcommand{\arraystretch}{1.25}
\begin{tabular}{ll}
    \textbf{s.t.}      & such that \\
    \textbf{coeff.}    & coefficient \\
    \textbf{iff}       & if and only if \\
    \textbf{rep.}      & representation \\
    \textbf{l.s.}      & linear system \\
    \textbf{l.i.}      & linearly independent \\
    \textbf{l.d.}      & linearly dependent \\
    \textbf{l.c.}      & linear combination \\
    \textbf{l.o.}      & linear operator \\
    \textbf{l.t.}      & linear transformation \\
    \textbf{v.s.}      & vector space \\
    \textbf{f.d.v.s.}  & finite dimensional vector space \\
    \textbf{n-d.v.s.}  & $n$-dimensional vector space \\
    \textbf{alt.}      & alternating \\
    \textbf{mult.}     & multilinear \\
\end{tabular}
\end{center}
\end{defbox}

% ------------------------------------------------------------
\begin{defbox}[Definition --- Number Sets and Field Notation]
The following symbols for standard number sets and fields are used throughout:
\begin{center}
\renewcommand{\arraystretch}{1.25}
\begin{tabular}{ll}
    $\N$                   & the set of natural numbers $\{1,2,3,\ldots\}$ \\
    $\N_0$                 & natural numbers including zero $\{0,1,2,3,\ldots\}$ \\
    $\Z$                   & the integers \\
    $\Q$                   & the rational numbers \\
    $\R$                   & the real numbers \\
    $\C$                   & the complex numbers \\
    $\F$                   & a general (arbitrary) field \\
    $\F_q$                 & a finite field with $q$ elements (Galois field) \\
    $\Z_p$ or $\F_p$       & the finite field of integers modulo a prime $p$ \\
\end{tabular}
\end{center}
\end{defbox}

% ------------------------------------------------------------
\begin{warnbox}[Notation Notice --- $\F$ vs.\ $k$ for a General Field]
In many algebra textbooks the letter $k$ is used for a general field.
\textbf{These notes use $\F = \mathbb{F}$ exclusively.}
Whenever you see $\F$, it denotes an arbitrary field unless stated otherwise.
\end{warnbox}

% ------------------------------------------------------------
\begin{tipbox}[Remark --- Additional Notation Introduced Later]
Further shorthand (e.g.\ $\Ker$, $\Img$, $\rk$, $\spn$, $\Hom$, $\End$,
$\Mmn{m}{n}$, $\Mn{n}$, $\chiA$, $\mA$, $\dime$) will be introduced and
defined at the point where the corresponding concepts first appear.
\end{tipbox}
